\section{MIFX Package Format}

A MIFX file is distributed as a compressed package containing the machining data and associated resources.

The package uses the file extension:

\begin{verbatim}
.mifx
\end{verbatim}

Internally, a MIFX file is a standard ZIP archive containing a structured directory tree.

\subsection{Root Structure}

The root of the package contains the main job descriptor and the primary directories used by the MIFX format.

Example structure:

\begin{verbatim}
job.mifx
|
+-- job.json
|
+-- setups/
+-- tools/
+-- geometry/
+-- in_process/
\-- addons/
\end{verbatim}

The only mandatory file at the root of the package is:

\begin{verbatim}
job.json
\end{verbatim}

This file acts as the entry point for interpreting the MIFX package.

\subsection{Directory Overview}

The standard directories of a MIFX package are summarized below.

\begin{itemize}

\item \textbf{setups/}

Contains machining setups and their associated operations.

\item \textbf{tools/}

Contains tool groups and tool definitions used by the machining process.

\item \textbf{geometry/}

Contains part, stock, fixture, and setup assembly geometry resources.

\item \textbf{in\_process/}

Contains optional intermediate workpiece states generated by simulation or analysis.

\item \textbf{addons/}

Contains optional extensions provided by external systems.

\end{itemize}

Implementations may omit directories that are not required by the machining job.

\subsection{Directory Flexibility}

The MIFX specification defines the logical structure of the package but does not require every directory to be present.

Systems generating MIFX packages may include only the resources necessary for the intended use.

Consumers of MIFX packages must ignore directories or files they do not recognize.

\subsection{Transport and Distribution}

Because the MIFX package is a ZIP archive, it can be easily transmitted, stored, or versioned using existing tools.

Examples include:

\begin{itemize}

\item email attachments
\item version control repositories
\item manufacturing data exchanges
\item digital job packages

\end{itemize}

This packaging approach ensures that all data required for interpreting the machining job can be distributed as a single file.